\sectioninfo{Introdução}{São Pio X}

\blettrine{U}{m} bom cristão, pela manhã, assim que desperta, deve fazer o sinal da Cruz, e oferecer o coração a Deus, dizendo estas ou outras palavras semelhantes: «Meu Deus, eu Vos dou o meu coração e a minha alma». Ao levantar da cama e enquanto nos vestimos, deveríamos pensar que Deus está presente, que este dia pode ser o último da nossa vida; ademais, devíamos levantar-nos e vestir-nos com toda a modéstia possível. A um bom cristão, apenas se tenha levantado e vestido, convém pôr-se na presença de Deus e ajoelhar-se, se pode, diante de alguma devota imagem, dizendo com devoção: «Eu Vos adoro, meu Deus, e Vos amo de todo o coração; dou-Vos graças por me terdes criado, feito cristão e conservado nesta noite; ofereço-Vos todas minhas acções, e peço-Vos que neste dia me preserveis do pecado, e me livreis de todo o mal. Assim seja». E rezar depois o Pai-Nosso, a Ave-Maria, o Credo, e os Actos de Fé, de Esperança e de Caridade, acompanhando-os com um vivo afecto do coração.\par
\begin{wrapfigure}{H}{0.5\textwidth}
\includegraphics[width=0.5\textwidth]{media/piox}
\end{wrapfigure}
O cristão, podendo, deveria todos os dias: 1º. Assistir com devoção à Santa Missa; 2º. Fazer uma visita, por breve que fosse, ao Santíssimo Sacramento; 3º. Rezar o terço do Santo Rosário.\par
Antes do trabalho, convém oferecê-lo a Deus, dizendo do coração: «Senhor, eu Vos ofereço este trabalho, dai-me a vossa bênção». Deve-se trabalhar para glória de Deus e para fazer a sua vontade. Antes da refeição, convém fazer o sinal da Cruz, estando de pé, e depois dizer com devoção: «Senhor, abençoai-nos a nós e ao alimento que vamos tomar, para nos conservarmos no vosso santo serviço». Depois da refeição, convém fazer o sinal da Cruz, e dizer: «Senhor, eu Vos dou graças pelo alimento que me destes; fazei-me digno de participar da mesa celeste».\par
Quando nos vemos atormentados por alguma tentação, devemos invocar com fé o Santíssimo Nome de Jesus ou de Maria, ou recitar fervorosamente alguma oração jaculatória, como, por exemplo: «Dai-me a graça, Senhor, de que eu nunca Vos ofenda»; ou então fazer o sinal da Cruz, evitando, porém, que as outras pessoas, pelos sinais externos, suspeitem da tentação. Quando uma pessoa reconhece ou receia ter cometido algum pecado, convém fazer imediatamente um acto de contrição, e procurar confessar-se quanto antes. Quando fora da igreja se ouve o sinal de elevação da Hóstia na Missa solene, ou da bênção do Santíssimo Sacramento é bom fazer, ao menos com o coração, um acto de adoração, dizendo, por exemplo: «Graças e louvores se dêem a todo o momento ao Santíssimo e diviníssimo Sacramento».\par
Ao toque das Ave-Marias pela manhã, ao meio-dia e à noite, o bom cristão recita o Anjo do Senhor «Angelus» com três Ave-Marias. À noite, antes de se deitar, convém pôr-se, como de manhã, na presença de Deus, recitar devotamente as mesmas orações, fazer um breve exame de consciência, e pedir perdão a Deus dos pecados cometidos durante o dia. Antes de adormecer, farei o sinal da Cruz, pensarei que posso morrer esta noite, e oferecerei o coração a Deus, dizendo: «Meu Senhor e meu Deus, eu Vos dou todo meu coração. Trindade Santíssima, concedei-me a graça de bem viver e de bem morrer. Jesus, Maria e José, eu Vos encomendo a minha alma».\par
No decurso do dia pode-se invocar a Deus frequentemente com as orações breves que se chamam «jaculatórias». Eis algumas:\par
«Senhor, valei-me»;\par
«Senhor, seja feita vossa santíssima vontade»;\par
«Meu Jesus, eu quero ser todo vosso»;\par
«Meu Jesus, misericórdia»;\par
«Doce Coração de Jesus, que tanto nos amou, fazei que eu Vos ame cada vez mais»;\par
«Doce Coração de Maria, sede minha salvação»;\par
É muito útil recitar, durante o dia, muitas jaculatórias, e podem recitar-se também com o coração, sem preferir palavras, caminhando, trabalhando, \&c. Além das orações jaculatórias, o cristão deveria exercitar-se na «mortificação cristã». Mortificar-se quer dizer privar-se, por amor a Deus, daquilo que agrada, e aceitar o que desagrada aos sentidos ou ao amor-próprio. Quando é o Santíssimo Sacramento levado a um enfermo, devemos, sendo possível, acompanhá-Lo com modéstia e recolhimento; e, se não é possível acompanhá-Lo, fazer um acto de adoração em qualquer lugar que nos encontremos, e dizer: «Consolai, Senhor, este enfermo, e concedei-lhe a graça de se conformar com vossa Santíssima vontade e de conseguir a sua salvação».\par
Ao tocar o sino pela agonia de algum moribundo, irei, se puder, à igreja orar por ele; e, não podendo, encomendarei a Nosso Senhor a sua alma, pensando que dentro em breve hei de encontrar-me também eu nesse estado. Ao ouvir sinais pela morte de alguém, procurarei rezar um «De profundis» ou um «Réquiem», ou um Pai-Nosso e uma Ave-Maria, pela alma desse defunto, e renovarei o pensamento da morte.
